
When a material permits spin-split bands while having compensated magnetic order, we speak of altermagnets.
In this thesis, altermagnets are being analyzed under the influence of different magnetic textures, breathing distortions of the lattice and Rashba-spin-orbit coupling (Rashba-SOC).  \\

For this, we introduce a two-dimensional tight binding model for the magnetic order of the T1 and T2 phase of \ce{Mn3Pt}, of \ce{Mn3Ge} and \ce{Mn3Sn}. An unitary transformation of the tight-binding Hamiltonian is used to achieve cell-periodic states. For the computation of the Berry curvature, we use the Fukui-Hatsugai-Suzuki method and additionally a composite representation for degenerate bands.
Subsequently, the band structure, Fermi-surface spin texture, Berry curvature and Chern number are computed for the different magnetic textures.\\
Here, we develop a theory for the symmetries holding on the magnetic textures and adapt them on the computations, confirming our theory. \\

In \ce{Mn3Ge} and \ce{Mn3Sn} we find a symmetry, imposing $C_{3z}$-symmetric behavior. The numerical computations show behavior, imposed by $\mathcal{TM}$ planes, for the T1 phase of \ce{Mn3Pt} and \ce{Mn3Ge}, which is assumed to hold in the presence of SOC.
In the nonbreathing limit without SOC, symmetries force the Berry curvature to vanish.
The T1 and T2 phase of \ce{Mn3Pt}, as well \ce{Mn3Ge} and \ce{Mn3Sn}, give equivalent results in the absence of SOC because their magnetic order differ by a global spin rotation.\\
Breathing distortions in the absence of SOC allow an odd Berry curvature in momentum, imposing a vanishing Chern number.\\
Using the received symmetries, we show that \ce{Mn3Pt} cannot exhibit an in-plane net magnetization, whereas \ce{Mn3Ge} \ce{Mn3Sn} can exhibit one along distinct directions in the presence of SOC. The examined materials remain spin-split in the absence of SOC, identifying them as strong altermagnets. Due to the in-plane net magnetization, \ce{Mn3Ge} and \ce{Mn3Sn} are identified as M-type altermagnets and T1 and T2 \ce{Mn3Pt} as S-type altermagnets.\\

% In general we find the symmetries $[E||C_{2z}]$, holding in the nonbreathing limit for  and $\Theta_s[C_{2z}||E]$ holding in the absence of SOC.

% \ce{Mn3Pt} is confirmed of S-type, where the T1 phase can exhibit an out-of plane net magnetization and T2 has a compensated net magnetization, as also stated in Ref.~\cite{Quelle3}.
% For \ce{Mn3Ge} and \ce{Mn3Sn} we find the $[C_{3z}||C_{3z}^{-1}]$-symmetry, enabling an in-plane net-magnetization along the $y$-axis, or along the $x$-axis, respectively, due to the additional $\mathcal{TM}$ and $\mathcal{M}$ planes.
% With the numerical computations, we assume the $\mathcal{TM}$ planes to be conserved in the presence of SOC.\\
% 
